\documentclass[11pt]{article}

% ─── Packages ─────────────────────────────────────────────────────
\usepackage[margin=1in]{geometry}
\usepackage{graphicx}
\usepackage{amsmath, amssymb}
\usepackage{booktabs}
\usepackage{hyperref}
\usepackage{xcolor}
\usepackage{enumitem}
\usepackage{caption}
\usepackage{subcaption}
\usepackage{float}
\usepackage{array}
\usepackage[numbers]{natbib}
\usepackage{tikz}
\usepackage{fancyhdr}
\usepackage{titlesec}
\usepackage{multirow}

\usetikzlibrary{arrows.meta, positioning, shapes.geometric, fit}

\graphicspath{{report_assets/report_assets/}}

\hypersetup{colorlinks=true, linkcolor=blue!60!black, citecolor=green!50!black, urlcolor=blue!70!black}

% ─── Header and Footer ─────────────────────────────────────────────
\pagestyle{fancy}
\fancyhf{}
\rfoot{\thepage}
\renewcommand{\headrulewidth}{0pt}

% ─── Title Formatting ─────────────────────────────────────────────
\title{\textbf{Classification of Indian Lepidoptera Species}\\
\Large Fine-Grained Butterfly Recognition Under Domain Shift}

\author{
  Akshit Bansal, Jayendra Singh, Kriti Chaturvedi,\\
  Sirjan Singh, Priyal Maheswari
}

\date{}

% ─── Section Formatting ─────────────────────────────────────────────
% \titleformat{\section}
%   {\normalfont\large\bfseries}
%   {}
%   {0em}
%   {}

% \titleformat{\subsection}
%   {\normalfont\normalsize\bfseries}
%   {}
%   {0em}
%   {}

\setlength{\parskip}{0.5em}
\setlength{\parindent}{0em}

\begin{document}
\maketitle

\section{Abstract}
Fine-grained classification of Indian Lepidoptera species presents compounding challenges: spatial information collapse in channel attention, progressive texture-cue loss through pooling, and extreme long-tail class imbalance. We conduct an 11-configuration ablation study across three experimental units--loss dynamics (30 epochs), architectural feature fusion (40 epochs), and transfer-learning layer freezing (40 epochs)--on a curated dataset of 24,825 images spanning 966 species. Our best configuration, cross-entropy with class-balanced sampling, achieves \textbf{87.06\% Top-1 accuracy} and \textbf{85.48\% macro-F1}, with balanced Dense/Sparse stratum performance (87.67\%/85.84\%). Coordinate Attention improves Sparse-class accuracy to 88.05\% by preserving positional wing-pattern information. Multi-Level Feature Interaction (MLFI) reaches 84.43\% Top-1 with strong Sparse performance (86.66\%), though its added parameterisation creates a 2-point gap versus the baseline. Head-only linear probing catastrophically fails at 53.20\%, confirming that ImageNet features are fundamentally misaligned with entomological morphology and full end-to-end fine-tuning is mandatory.

\section{Introduction and Motivation}

India's 1,500+ Lepidoptera species represent one of the world's richest butterfly biodiversity hotspots. Automated classification from field photographs is critical for ecological monitoring and conservation planning. Yet this task fundamentally differs from standard benchmarks: discriminative features are spatially precise (eyespot location separates genera), hierarchically distributed across scales (texture bands vs overall wing morphology), and unevenly sampled (dominant species exceed 500 images while 400 rare species have $<$50 images).

Standard CNNs fail on this task in predictable ways. Channel attention mechanisms like SE-Net~\cite{hu2018squeeze} collapse spatial dimensions through global average pooling, discarding the positional information essential for wing pattern discrimination. Deep pooling progressively destroys fine-grained texture cues, routing only high-level semantics to the classifier. Uniform sampling over-indexes on common species, starving rare classes of gradient signal. We structure our investigation around three questions: (1)~how do loss formulations and sampling strategies interact under long-tail distributions, (2)~where is the optimal point for integrating attention and multi-level fusion, and (3)~what is the optimal adaptation pathway when transferring from non-biological source domains.

\section{Dataset and Experimental Setup}

The IndoLepAtlas dataset comprises 24,825 images (966 species post-filtering) curated from the Indian Lepidoptera web atlas and iNaturalist, with rigorous metadata: GPS coordinates, observation dates, and 9-zone biogeographic mapping across 34 Indian states. Geographic contributors: Maharashtra (4,383), Karnataka (3,065), Arunachal Pradesh (2,918). Temporal coverage spans all 12 months with monsoon peak (Aug--Nov). Class distribution ranges from 500+ images (dominant species like \emph{Papilio polytes}) to $<$5 images (tail species), with key metadata coverage of 98.3\% (location) and 76.1\% (date). Classes are stratified into \textbf{Dense} ($\geq$50 images) and \textbf{Sparse} ($<$50 images) strata for evaluation.

We execute a systematic 11-configuration ablation across three experimental units on a DGX V100 cluster. Unit~I (loss dynamics) trains for 30 epochs; Units~II and III (architecture and freezing) train for 40 epochs. All experiments use a ConvNeXt-Tiny backbone~\cite{liu2022convnext} with AdamW optimisation (differential learning rates: backbone $0.1\times$, new heads $10^{-4}$), cosine annealing with 5-epoch warmup, mixed precision AMP, and macro-F1 checkpoint selection to ensure long-tail sensitivity.

\section{Core Findings}

\subsection{Unit I: Long-Tail Loss Dynamics (30 Epochs)}

\textbf{Question:} \emph{How do different loss formulations and sampling strategies impact representational capacity on long-tailed fine-grained distributions?}

We test cross-entropy (CE) and Focal Loss~\cite{lin2017focal} across balanced and unbalanced sampling. Results are summarised in Table~\ref{tab:master} (Unit~I rows).

\begin{figure}[H]
\centering
\includegraphics[width=0.95\linewidth]{plots/unit1_training_curves.png}
\caption{Unit~I training curves: validation loss and accuracy over 30 epochs for all four loss--sampling configurations.}
\label{fig:unit1_curves}
\end{figure}

\begin{figure}[H]
\centering
\includegraphics[width=0.85\linewidth]{plots/unit1_bar_comparison.png}
\caption{Unit~I metric comparison: Top-1 accuracy, Macro-F1, Dense and Sparse accuracy across loss--sampling strategies.}
\label{fig:unit1_bars}
\end{figure}

\textbf{The ``Unbalanced'' illusion.} CE~+~Unbalanced achieves the highest raw Top-1 accuracy (87.60\%), but this masks severe stratum imbalance: Dense accuracy inflates to 91.22\% while Sparse accuracy drops to 80.35\%. The macro-F1 plummets to 80.23\%--a 5.25-point gap below the balanced variant--confirming the model abandons rare species to maximise majority-class performance.

\textbf{CE + Balanced as optimal compromise.} Balanced sampling forces equitable gradient allocation, compressing the Dense--Sparse gap to just 1.83 points (87.67\% vs 85.84\%) while achieving the highest macro-F1 (85.48\%) and macro-precision (87.43\%). The mechanism is direct: inverse-frequency sampling ensures rare-class gradients remain proportional regardless of natural distribution skew.

\textbf{The Focal Loss penalty.} Focal Loss~\cite{lin2017focal} with $\gamma=2$ achieves the highest Sparse-class accuracy (86.40\% under balanced sampling) but severely degrades Dense-class performance to 80.85\%. In 966-class fine-grained settings, the $(1-p_t)^\gamma$ modulation over-penalises well-classified easy examples (typically Dense strata), damaging foundational feature extraction for common species. Focal~+~Unbalanced partially recovers Dense accuracy (85.70\%) but at the cost of Sparse performance. This contradicts the assumption that loss-function sophistication universally solves long-tail problems--explicit sampling proves more reliable.

\subsection{Unit II: Feature-Fusion Ablation (40 Epochs)}

\textbf{Question:} \emph{Where is the optimal structural transition point for integrating attention mechanisms, and does multi-level feature integration improve fine-grained classification?}

We ablate structural complexity across three progressive phases (Table~\ref{tab:master}, Unit~II rows).

\begin{figure}[H]
\centering
\includegraphics[width=0.95\linewidth]{plots/unit2_training_curves.png}
\caption{Unit~II training curves: validation loss and accuracy over 40 epochs for Baseline, +CA, and +CA+MLFI architectures.}
\label{fig:unit2_curves}
\end{figure}

\begin{figure}[H]
\centering
\includegraphics[width=0.85\linewidth]{plots/unit2_bar_comparison.png}
\caption{Unit~II metric comparison: progressive feature-fusion architecture performance.}
\label{fig:unit2_bars}
\end{figure}

\textbf{Coordinate Attention preserves spatial morphology.} Introducing Coordinate Attention~\cite{hou2021coordinate} after each backbone stage (Phase~2) yields the best overall configuration: 86.77\% Top-1, 85.75\% macro-F1, and a notable 1.29-point improvement in Sparse-class accuracy (86.76\% $\to$ 88.05\%). CA decomposes channel-spatial attention into two 1D processes: $\mathbf{z}_h = \text{AvgPool}_W(\mathbf{X})$ and $\mathbf{z}_w = \text{AvgPool}_H(\mathbf{X})$, preserving directional structure naturally suited to butterfly morphology (horizontal banding, vertical vein streaking). Unlike SE-Net's global average pooling, CA retains \emph{where} a pattern occurs--critical when eyespot position on forewing vs hindwing separates genera. The overhead is minimal ($<$0.5\% added parameters), making CA a high-value, low-cost intervention.

\textbf{MLFI: powerful but over-parameterised.} Phase~3 concatenates features from all four backbone stages (192 + 384 + 96 + 768 = 1,440 dimensions)~\cite{yuan2025fgbnet} to preserve texture cues destroyed by successive pooling. This adds approximately 12M parameters. At 40 epochs, MLFI reaches 84.43\% Top-1 and 86.66\% Sparse accuracy--demonstrating that multi-level fusion does benefit rare classes--but creates a 2.07-point Top-1 gap versus the baseline. The newly initialised MLFI fusion layers (Xavier initialisation, full learning rate) create an optimisation imbalance with the near-converged backbone ($0.1\times$ learning rate), resulting in a poorly conditioned loss landscape. Because early layers capture high-frequency textural variance that varies fiercely under different lighting, forcing these raw activations into the classification head introduces feature noise that overwhelms the refined late-layer semantic signals.

\subsection{Unit III: Transfer Learning Layer Freezing (40 Epochs)}

\textbf{Question:} \emph{Which hierarchical layers contain domain-agnostic representations versus domain-specific concepts, and what is the optimal adaptation pathway for non-standard image domains?}

\begin{figure}[H]
\centering
\includegraphics[width=0.95\linewidth]{plots/unit3_training_curves.png}
\caption{Unit~III training curves: validation loss and accuracy over 40 epochs for all four layer-freezing strategies.}
\label{fig:unit3_curves}
\end{figure}

\begin{figure}[H]
\centering
\includegraphics[width=0.85\linewidth]{plots/unit3_bar_comparison.png}
\caption{Unit~III metric comparison: layer-freezing strategy performance (see Table~\ref{tab:master}, Unit~III rows).}
\label{fig:unit3_bars}
\end{figure}

\textbf{The domain-shift wall.} Head-only linear probing--freezing the entire backbone and training only a classification head--achieves just 53.20\% Top-1 on 966 species. Dense-class accuracy collapses to 44.46\%, worse than Sparse-class accuracy (70.75\%), a striking inversion that reveals ImageNet features are not merely \emph{incomplete} but actively \emph{misaligned} with Lepidoptera morphology. SE-Net activations encode ``furry surface textures'' and ``rounded shapes,'' not ``submarginal band'' or ``discal spot.'' Fine-grained biological classification is fundamentally a domain-shift problem~\cite{dosovitski2021image} disguised as a classification problem.

\textbf{Early vs late block freezing.} Freezing late blocks (84.33\%) slightly outperforms freezing early blocks (83.77\%). This implies that the early-layer Gabor-like edge detectors require significant recalibration to adapt to butterfly wing scales, biological venation, and field photography conditions. Adapting early layers creates a stronger foundation for the frozen late layers to build upon. The inter-strategy variance is small ($\Delta < 1\%$), indicating that under sufficient training budget, any partial-freezing strategy approaches end-to-end performance.

\textbf{End-to-end supremacy.} Full gradient backpropagation achieves the best macro-F1 (84.11\%) and the most balanced stratum performance (Dense 83.89\%, Sparse 86.61\%). Because the semantic gap between ImageNet (everyday objects) and IndoLepAtlas (entomological morphology) is profound, the entire representational manifold must be warped for this domain.

\section{Cross-Unit Analysis}

\subsection{Stratum-Level Performance}

A consistent pattern emerges across all 11 configurations: Sparse-class accuracy often exceeds Dense-class accuracy under balanced training (Table~\ref{tab:master}). This occurs because balanced sampling provides proportionally more gradient signal per Sparse-class sample, while Dense classes--despite more data--contain greater intra-class morphological variation (seasonal dimorphism, sexual dimorphism, geographic colour variation). The most pronounced effect appears in Phase~2 (CA), where Sparse accuracy reaches 88.05\% versus Dense accuracy of 86.13\%, suggesting that Coordinate Attention's spatial preservation is especially valuable for classes with few but morphologically distinctive exemplars.

\begin{figure}[H]
\centering
\includegraphics[width=0.85\linewidth]{plots/stratum_comparison.png}
\caption{Dense vs Sparse accuracy across all 11 configurations, highlighting the stratum-equalising effect of balanced sampling and attention mechanisms.}
\label{fig:stratum}
\end{figure}

\subsection{Top-5 Accuracy Saturation}

All configurations except Head-Only achieve Top-5 accuracy above 96.6\%, indicating the model's confusion is largely confined to the most morphologically similar species pairs. This suggests hierarchical classification strategies or pairwise discriminators as a high-leverage future direction.

\begin{figure}[H]
\centering
\includegraphics[width=0.75\linewidth]{plots/best_models_radar.png}
\caption{Radar chart comparing the best configuration from each unit across five evaluation axes.}
\label{fig:radar}
\end{figure}

\begin{table}[H]
\centering
\caption{Complete 11-configuration ablation summary.}
\label{tab:master}
\small
\begin{tabular}{@{}llccccc@{}}
\toprule
\textbf{Unit} & \textbf{Configuration} & \textbf{Epochs} & \textbf{Top-1} & \textbf{Macro-F1} & \textbf{Dense} & \textbf{Sparse} \\
\midrule
\multirow{4}{*}{I}
  & CE + Balanced       & 30 & \textbf{87.06} & \textbf{85.48} & 87.67 & 85.84 \\
  & CE + Unbalanced     & 30 & 87.60 & 80.23 & \textbf{91.22} & 80.35 \\
  & Focal + Balanced    & 30 & 82.69 & 82.64 & 80.85 & \textbf{86.40} \\
  & Focal + Unbalanced  & 30 & 84.83 & 82.70 & 85.70 & 83.07 \\
\midrule
\multirow{3}{*}{II}
  & Phase 1 (Baseline)    & 40 & 86.50 & 85.65 & \textbf{86.36} & 86.76 \\
  & Phase 2 (+ CA)        & 40 & \textbf{86.77} & \textbf{85.75} & 86.13 & \textbf{88.05} \\
  & Phase 3 (+ CA + MLFI) & 40 & 84.43 & 83.34 & 83.32 & 86.66 \\
\midrule
\multirow{4}{*}{III}
  & End-to-End      & 40 & \textbf{84.79} & \textbf{84.11} & \textbf{83.89} & 86.61 \\
  & Head Only       & 40 & 53.20 & 57.02 & 44.46 & 70.75 \\
  & Freeze Early    & 40 & 83.77 & 83.20 & 82.74 & 85.84 \\
  & Freeze Late     & 40 & 84.33 & 84.02 & 83.15 & \textbf{86.71} \\
\bottomrule
\end{tabular}
\end{table}

\section{Discussion}

Three structural constraints drive the failure modes observed across our ablation:
\begin{enumerate}[label=(\arabic*)]
    \item \textbf{Spatial collapse} in standard attention mechanisms obscures the precise localisation needed for wing pattern discrimination. Coordinate Attention addresses this by preserving directional spatial maps, yielding the highest Sparse-class accuracy (88.05\%) across all configurations.
    \item \textbf{Progressive texture loss} through repeated pooling destroys fine-grained cues essential for species differentiation. MLFI addresses this through early-layer feature preservation, achieving strong Sparse performance (86.66\%), but the added parameterisation creates optimisation difficulties that limit overall accuracy.
    \item \textbf{Class imbalance} under uniform sampling starves rare species of gradient signal. Balanced CE sampling achieves the highest macro-F1 (85.48\%) by compressing the Dense--Sparse accuracy gap to 1.83 points.
\end{enumerate}

The domain-shift finding is the most generalisable insight. The Head-Only configuration's 53.20\% accuracy--with Dense-class accuracy (44.46\%) falling below Sparse-class accuracy (70.75\%)-- demonstrates that ImageNet features are not merely insufficient but structurally inappropriate for entomological morphology. The Dense-class inversion occurs because ImageNet's learned features (object silhouettes, fur textures, scene context) actively interfere with distinguishing morphologically similar dominant species, while Sparse-class species often possess more distinctive visual signatures that survive even frozen-feature extraction.

These findings extend beyond butterflies to medical imaging, satellite-based crop monitoring, and microscopy-based cell classification--all domains facing similar constraints of extreme class imbalance, domain-specific morphological features, and limited per-class training data~\cite{amarathunga2022thrips}.

\section{Conclusion}

Three key findings emerge from our 11-configuration ablation:
\begin{enumerate}[label=(\arabic*)]
    \item \textbf{Balanced sampling dominates:} CE + Balanced achieves the best macro-F1 (85.48\%) with the most equitable stratum performance. Focal Loss improves Sparse-class accuracy but at the cost of Dense-class performance.
    \item \textbf{Coordinate Attention is high-value, low-cost:} CA adds minimal parameters yet improves Sparse-class accuracy by 1.29 points (88.05\%), making it the optimal architectural intervention. MLFI fusion shows promise for rare classes but requires further optimisation to close the 2-point Top-1 gap.
    \item \textbf{ImageNet transfer fails catastrophically:} Head-Only achieves 53.20\%--full end-to-end fine-tuning is mandatory for specialised biological domains.
\end{enumerate}

Future work includes geotemporal fusion using seasonal and biogeographic embeddings, prototypical network heads~\cite{snell2017prototypical} for the $\sim$400 classes with $<$50 images, and domain-specific self-supervised pretraining (DINO on butterfly images)~\cite{feichtenhofer2022masked} to bootstrap features from unlabelled in-domain data.

\newpage

\bibliographystyle{plainnat}

\begin{thebibliography}{9}

\bibitem[Yuan et~al., 2025]{yuan2025fgbnet}
Y.~Yuan et~al., FGBNet: A Bio-Subspecies Classification Network with Multi-Level Feature Interaction, \emph{Diversity}, 2025. \url{https://doi.org/10.3390/d17040237}

\bibitem[Hou et~al., 2021]{hou2021coordinate}
Q.~Hou, D.~Zhou, and J.~Feng, Coordinate Attention for Efficient Mobile Network Design, \emph{CVPR}, 2021. \url{https://arxiv.org/abs/2103.02907}

\bibitem[Liu et~al., 2022]{liu2022convnext}
Z.~Liu, H.~Mao, C.-Y.~Wu, C.~Feichtenholer, T.~Darrell, and S.~Xie, A ConvNet for the 2020s, \emph{CVPR}, 2022. \url{https://arxiv.org/abs/2201.03545}

\bibitem[Hu et~al., 2018]{hu2018squeeze}
J.~Hu, L.~Shen, and G.~Sun, Squeeze-and-Excitation Networks, \emph{CVPR}, 2018. \url{https://arxiv.org/abs/1709.01507}

\bibitem[Lin et~al., 2017]{lin2017focal}
T.-Y.~Lin, P.~Goyal, R.~Girshick, K.~He, and P.~Doll\'{a}r, Focal Loss for Dense Object Detection, \emph{ICCV}, 2017. \url{https://arxiv.org/abs/1708.02002}

\bibitem[Amarathunga et~al., 2022]{amarathunga2022thrips}
D.~C.~Amarathunga et~al., Fine-grained image classification of microscopic insect pest species, \emph{Computers and Electronics in Agriculture}, 2022. \url{https://doi.org/10.1016/j.compag.2022.107462}

\bibitem[Snell et~al., 2017]{snell2017prototypical}
J.~Snell, K.~Swersky, and R.~Zemel, Prototypical Networks for Few-Shot Learning, \emph{NeurIPS}, 2017. \url{https://arxiv.org/abs/1703.05175}

\bibitem[Feichtenhofer et~al., 2022]{feichtenhofer2022masked}
C.~Feichtenhofer, H.~Fan, Y.~Li, and K.~He, Masked autoencoders as spatiotemporal learners, \emph{NeurIPS}, 2022. \url{https://proceedings.neurips.cc/paper_files/paper/2022/file/e97d1081481a4017df96b51be31001d3-Paper-Conference.pdf}

\bibitem[Dosovitskiy et~al., 2021]{dosovitski2021image}
A.~Dosovitskiy et~al., An image is worth 16x16 words: Transformers for image recognition at scale, \emph{ICLR}, 2021. \url{https://arxiv.org/pdf/2010.11929}

\bibitem[Loshchilov and Hutter, 2017]{loshchilov2017sgdr}
I.~Loshchilov and F.~Hutter, SGDR: Stochastic Gradient Descent with Warm Restarts, \emph{ICLR}, 2017. \url{https://arxiv.org/abs/1608.03983}

\end{thebibliography}

\end{document}